\chapter*{Abstract}
\addcontentsline{toc}{chapter}{Abstract}

% Generative Artificial Intelligence (\acrshort{ai}) systems have enabled the large scale creation of images, with these technologies being available to the public. However, these systems creation is largely dependent on the existence of large web-scraped datasets serving as training data such as Re-LAION-5B serving as the focus of this research alongside the Stable Diffusion v1.5 model trained on it. Due to the nature of the training data, their has been increased concern regarding the ethical implications imposed particularly related to Bias and unintended correlations embedded within the training data. Particularly non-semantic or spurious features present within the training dataset may serve as shortcuts for sensitive demographic attributes, negatively impacting the downstream model behaviour and fairness outcomes.

Generative Artificial Intelligence (AI) systems have enabled the large-scale creation of images, with these technologies becoming increasingly accessible to the public. However, the development of these systems is largely dependent on the existence of massive web-scraped datasets serving as training data, such as Re-LAION-5B, which forms the focus of this research alongside the Stable Diffusion v1.5 model trained on it. Due to the nature of web-scraped training data, there has been increased concern regarding the ethical implications imposed, particularly related to bias and unintended correlations embedded within the training corpus. Specifically, non-semantic or spurious features present within the training dataset may serve as shortcuts for various downstream tasks, negatively impacting model behaviour and fairness outcomes. While these shortcuts can emerge for any task, demographic attributes represent a particularly concerning application, alongside other attributes unrelated to the inherent semantic content of the dataset.

% This thesis investigates the demographic bias within the Re-LAION-5B dataset and resultant generative images via the SDv1.5 model with a direct focus on identifying spurious features within these two aspects as well as identifying if the model amplifies, alters or mitigated those spurious features identified within the respective training corpus. The study adopts a classifier-based approach derived from prior research, in which models are trained to distinguish between source datasets utilising images which have underwent specific transformations to reduce the information available to the model and force the reliance on specific image features, thereby exposing the presence of dataset specific spurious features.

This thesis investigates the demographic bias within the Re-LAION-5B dataset and resultant generative images produced by the Stable Diffusion v1.5 model, with a direct focus on identifying spurious features within these two components and determining whether the generative model amplifies, alters, or mitigates those spurious features identified within the respective training corpus. The study adopts a classifier-based approach derived from prior research, in which discriminative models are trained to distinguish between Re-LAION-5B and Stable Diffusion v1.5 generated images. These classifiers utilise images that have undergone specific transformations. These transformations systematically reduce the semantic information available to the model and force reliance on specific low-level image features, thereby exposing the presence of dataset specific spurious patterns that can be exploited by the model for image source identification.

% Building upon this analysis, the relevant image sets are further annotated in terms of gender and skin-tone in order to further investigate the dataset specific spurious features allowing for the detection of harmful features serving as model shortcuts for deriving these sensitive attributes. To further contextualise these findings, the demographic distributions of retrieved and generated images across 200 occupation categories are compared against real-world statistics.

Building upon this classifier based analysis, the relevant image sets used in the discrimination task are further annotated in terms of gender and skin-tone attributes. This secondary analysis investigates whether the previously identified dataset specific spurious features demonstrate correlations with these sensitive demographic attributes, allowing for the detection of harmful features that may serve as model shortcuts for deriving gender or skin-tone classifications. To further contextualise these findings and assess the real world implications of observed biases, the demographic distributions of retrieved and generated images across 200 occupation categories are compared against real world demographic statistics, revealing patterns of representation bias and potential amplification effects in the generative model.

% Finally, in line with the spurious features identified in the prior sections, the thesis outlines a set of inference-time strategises for mitigating these biases in Stable Diffusion v1.5 generated image without requiring full model retraining. By linking dataset characteristics, spurious feature learning, and generative bias, this work contributes practical insights into the evaluation of large-scale multimodal datasets and the development of fairer text-to-image generation systems.

Finally, based on the spurious features identified in the prior sections, this thesis outlines a comprehensive set of inference-time strategies for debiasing the Stable Diffusion v1.5 model without requiring full model retraining. These proposed methods aim to mitigate the effect that identified spurious features have on resultant generated images, offering practical approaches for improving fairness in deployed systems. By explicitly linking dataset characteristics, spurious feature learning mechanisms, and generative bias outcomes, this work contributes both methodological insights for evaluating large-scale multimodal datasets and practical interventions for the development of fairer text-to-image generation systems.
